\documentclass[11pt]{article}
\usepackage[T1]{fontenc}
\usepackage{textcomp}
\usepackage[margin=0.75in]{geometry}
\usepackage{amsmath,amssymb,amsthm}
\usepackage{array,booktabs}
\usepackage{hyperref}
\setlength{\parindent}{0pt}
\newtheorem{theorem}{Theorem}
\theoremstyle{definition}
\newtheorem{definition}{Definition}
\title{Flow Proof Artifact: prop-15-reciprocal-sides-implies-equal}
\author{Generated by \texttt{flow doc proof}}
\date{Flow Proof Book}
\begin{document}
\maketitle
Equiangular parallelograms in which the sides about the equal angles are reciprocally proportional are equal.\\[0.5em]
\textbf{Source.} euclid — Elements, Book VI, Proposition 15\\[0.5em]
\bigskip
\noindent{\large \textbf{Derived fact 1.} \textit{Equiangular parallelograms in which the sides about the equal angles are reciprocally proportional are equal} \hfill the Euclidean plane \cdot Euclid Book VI \cdot Proposition 15: equiangular parallelograms with reciprocal sides are equal}\par
\smallskip\noindent\textit{Coordinate: the Euclidean plane \cdot Euclid Book VI \cdot Proposition 15: equiangular parallelograms with reciprocal sides are equal}\par
\smallskip\noindent\textit{Source: euclid — Elements, Book VI, Proposition 15}\par
\smallskip\noindent\textit{Built on: proposition 14: equal equiangular parallelograms have reciprocally proportional sides, for Euclid Book VI on the Euclidean plane}\par
\medskip
\noindent\textbf{Goal.} Equiangular parallelograms in which the sides about the equal angles are reciprocally proportional are equal\par
\noindent$\displaystyle \text{reciprocal sides implies equal}$\par
\medskip
\noindent
\begin{tabular}{@{} >{\raggedright\arraybackslash}p{0.06\textwidth} >{\raggedright\arraybackslash}p{0.47\textwidth} >{\raggedleft\arraybackslash}p{0.41\textwidth} @{}}
\textbf{\#} & \textbf{Proof} & \textbf{Mathematics} \\
\midrule
\textbf{1.} & We prove that proposition 15: equiangular parallelograms with reciprocal sides are equal for Euclid Book VI the Euclidean plane. & \\[0.45em]
\textbf{2.} & Let ABCD = parallelogram\_1. & \\[0.45em]
\textbf{3.} & Let EFGH = parallelogram\_2. & \\[0.45em]
\textbf{4.} & Let reciprocal = side\_ratio. & \\[0.45em]
\textbf{5.} & From step 2, step 3, and step 4, we can deduce that area ABCD equals area EFGH. & $\displaystyle \text{area ABCD} = \text{area EFGH}$ \\[0.45em]
\textbf{6.} & We invoke the derived fact governing Euclid Book VI the Euclidean plane: proposition 14: equal equiangular parallelograms have reciprocally proportional sides, for Euclid Book VI on the Euclidean plane. & \\[0.45em]
\textbf{7.} & From step 6, this implies reciprocal sides implies equal. Hence proven. & $\displaystyle \text{reciprocal sides implies equal}$ \\[0.45em]
\bottomrule
\end{tabular}
\label{thm:Geometry-Euclid Book VI-proposition_15_equiangular_parallelograms_with_reciprocal_sides_are_equal}
\par\medskip\hrule
\end{document}
